\section{Cross-Language and Cross-Structure Equivalence}
\label{sec:equivalence}

Speed differences are only meaningful if the implementations compute the same
thing. Two equivalence gates run before any timing is reported, and both
commit their evidence next to the results.

\subsection{Gate 1: the four languages agree}
\label{sec:equivalence-languages}

The checker (\texttt{src/utils/benchmarks/equivalence\_check.py}) selects
\fixed{100} books from the corpus, evenly spaced by ID, and runs every
port's \texttt{equivalence} mode over them. Each port writes one canonical
file with three sections:

\begin{enumerate}
    \item \texttt{tokens} --- the token stream of every sample book
    (\texttt{tokens<TAB>book\_id<TAB>token...});
    \item \texttt{postings} --- the inverted index built from those tokens
    (\texttt{postings<TAB>term<TAB>id...});
    \item \texttt{query} --- the answer to every one of the \fixed{30}
    shared queries (\texttt{query<TAB>ordinal<TAB>id...}).
\end{enumerate}

The four files are compared byte for byte. The run recorded in
\texttt{benchmarks\_results/equivalence\_report.json} produced the same
SHA-256 digest in all four languages,

\begin{center}
    \texttt{cdaa9369cbd91c4f89550efdcdd7c5c21e61adc97bd49f6e415b8e57d65ce2e4},
\end{center}

with \texttt{all\_equal: true}. Identical bytes imply identical tokens,
identical postings and identical query answers: the cross-language latency
comparison in \cref{sec:language-performance} compares the same work
everywhere. The tournament refuses to run any timed phase if this gate fails.

\subsection{Gate 2: the three structures agree}
\label{sec:equivalence-structures}

Within a language, the latency chart compares three index structures. The
test \texttt{tests/test\_datamarts/test\_structure\_equivalence.py} builds
the same corpus into the JSON file, the folder index and MongoDB, then
queries \emph{every} term of the corpus plus absent terms on each structure
and asserts that all answers are identical. A second case repeats the check
\emph{after} adding the same new book to each structure through its own
incremental-update path, so an update cannot make one structure diverge
silently. All cases pass (\cref{tab:equivalence-tests}); the MongoDB cases
skip themselves with an explicit reason when the server is down, and the
tournament only runs timed MongoDB phases when the server is reachable.

\begin{table}[htbp]
    \centering
    \caption{Cross-structure equivalence cases and results.}
    \label{tab:equivalence-tests}
    \begin{tabular}{@{}llc@{}}
        \toprule
        Case & Structures compared & Result \\
        \midrule
        All terms + absent terms, fresh build & JSON = folder & identical \\
        All terms + absent terms, fresh build & JSON = folder = MongoDB & identical \\
        Re-queried after the same incremental update & JSON = folder = MongoDB & identical \\
        \bottomrule
    \end{tabular}
\end{table}

\subsection{What equivalence does and does not prove}
\label{sec:equivalence-scope}

The gates prove \emph{functional} equivalence on this corpus and workload:
same tokens, same postings, same answers. They say nothing about speed ---
that is exactly what the following sections measure. Conversely, the
latency charts inherit the gates' guarantee: every bar in them comes from an
implementation that is known to answer identically.
